

\ksec{Class Semantics (KerML 8.4.4.3)}
\pn \emph{Classes}\index{class} are classifiers that classify \emph{occurrences}\index{occurrence}, which exist in time and space (see KerML \S9.2.4). Relations between an occurrence and other things can change over time and space, while the occurrence still maintains an independent identity.

\pn A class is declared as a classifier (see KerML \S7.3.3), using the keyword \lstinline{class}. If no owned superclassing is explicitly given for the class, then it is implicitly given a default superclassing to the class \texttt{Occurrence} from the \texttt{Occurrences} model library (see KerML \S9.2.4).

%\pn Therefore, temporal semantics is determined in terms of Core semantics by annotating features of \texttt{Occurrences\-::Occurrence} with \texttt{@Assert}\index{@Assert} metadata.
%The temporal meaning of features of \texttt{Occurrence} are defined by semantic metadata within the body, \{ \}, of whatever to which the \texttt{@Assert} applies.
%
%\begin{lstlisting}
%metaclass Assert specializes SemanticMetadata {
%	doc
%	/*
%	 * Assert attaches temporal logic formulas to KerML elements
%	 */
%	feature n[0..1] : String;  //name
%	feature f[1..*] : String;  //well-formed formula(s)
%	feature t[0..*] : String;  //theorem name(s) which attests to this @Assert
%}
%\end{lstlisting}

\subsection{Occurrence Interval}
\pn As discussed before, temporal logics use either instant or intervals as their fundamental element.
Necessarily, the modal logic uses instants, so \texttt{Occurrence} will be an interval defined by instants.
\comment{This is the biggest and most controversial change.}

\pn  \texttt{Occurrence} is a \texttt{Domain::Interval} having \texttt{startShot} as  \texttt{lowerBound}, and  \texttt{endShot} as  \texttt{upperBound}.
\texttt{Occurrence} is default closed on both ends.
  \texttt{middleTimeSlice} is the open interval between \texttt{startShot} and \texttt{endShot} when they are unequal.

\ecaption{Occurrence}  
\begin{lstlisting}
abstract class Occurrence specializes Anything, Domain::Interval 
    disjoint from DataValue {
  feature :>> openLeft default false { }
  feature :>> openRight default false { }
  portion feature startShot: Time::Instant[1] :>> lowerBound { }
  portion feature endShot: Time::Instant[1] :>> upperBound {}
  portion feature middleTimeSlice: Interval[0..1] {
    language "Domain" /* middleTimeSlice = startShot ,, endShot" */
    :>> lowerBound = startShot;
    >> upperBound = endShot;
    :>> openLeft = true;
    :>> openRight = true;
  }
  inv { isEmpty((that as Occurrence).middleTimeSlice) == 
    ((that as Occurrence).startShot == (that as Occurrence).endShot) }
\end{lstlisting}

\subsection{Replacing M0 Features}\label{subsec:replacingM0}

\pn Precise Semantics of UML State Machines (PSSM) is an OMG standard \cite{PSSM}.
It defines behavior of state machines, which are \comment{hideously-complex} UML data structures, by constructing other 
\comment{even more hideously complex} UML data structures which change over time.  State machine models (M1) are distinct from their executions (M0).

\pn Unfortunately, standard libraries for KerML and SysML prior to SFS mix M0 execution with M1 model, starting with \texttt{Occurrence}.
Features like \texttt{predecessors} are lists of \texttt{Occurrence}, in this case every other  \texttt{Occurrence} prior to this one. 
\texttt{predecessors} does not represent the (static) model M1, but rather execution M0.  

\pn Mixing M0 features with M1 features caused confusion; they are not distinguished.  Furthermore, no execution, even in simulation, will construct a conforming UML data structure.  M0 features are purely notional.

\pn Worse, M0 features like \texttt{successors} implicitly assume eternalism metaphysics (\ref{sec:fortime}).   \texttt{successors} cannot have values until execution is finished.  Post hoc semantics does not define what the system should do \emph{now}.

\pn Fortunately, M0 features are overwhelmingly used to test membership.  An \texttt{Occurrence} happens before this one if it is one of the \texttt{predecessors}.  The valuable concepts expressed by M0 features can be expressed by functions and predicates instead.\footnote{
This presumes for SFS that expressions, functions, and predicates just determine values, and are not \texttt{Occurrence} themselves (\ref{}).}

\pn Use function \texttt{getLife} instead of feature \texttt{portionOfLife}.

\ecaption{getLife}  
\begin{lstlisting}
		function getLife {
			doc /* Life is top-level Occurrence, suboccurrence of no other */
			language "Domain" /* during(this,result) and
			   not exists o : Occurrence in o<>this that during(o,result) */	
			return result : Life[1];
			if superoccurrence->isEmpty() ? this as Life
			  else superoccurrence->head().getLife	
		}
\end{lstlisting}

\pn Use predicate \texttt{sameLife} instead of feature \texttt{sameLifeOccurrences}.

\ecaption{sameLife}  
\begin{lstlisting}
		predicate sameLife {
			doc /* SFS: o has the same Life as this */
			in o : Occurrence;
			this.getLife == o.getLife
		}
\end{lstlisting}

\pn Use predicate \texttt{withoutOccurrence} instead of feature \texttt{withoutOccurrences}.

\ecaption{withoutOccurrence}  
\begin{lstlisting}
    predicate withoutOccurrence {  
      doc /* SFS: true when o is separated in either time or space */
      in o : Occurrence;
      nonoverlaps(o,this) or not regionOverlap(location(o),location(this)) 
    }
\end{lstlisting}

\pn Use predicate \texttt{predecessor} instead of feature \texttt{predecessors}.

\ecaption{predecessor}  
\begin{lstlisting}
    predicate predecessor { 
      doc /* SFS:  true when o precedes this */
      in o : Occurrence;
      precedes(o,this)
    }
\end{lstlisting}

\pn Use predicate \texttt{successor} instead of feature \texttt{successors}.

\ecaption{successor}  
\begin{lstlisting}
    predicate successor {
      doc /* SFS:  true when self precedes o */
      in o : Occurrence;
      precedes(this,o)
    }
\end{lstlisting}

\pn Use predicate \texttt{immediatePredecessor} instead of feature \texttt{immediatePredecessors}.

\ecaption{immediatePredecessor}  
\begin{lstlisting}
    predicate immediatePredecessor { 
      doc /* SFS:  true when o nearly meets this */
      in o : Occurrence;
      nearlyMeets(o,this)
    }
\end{lstlisting}

\pn Use predicate \texttt{immediateSuccessor} instead of feature \texttt{immediateSuccessors}.

\ecaption{immediateSuccessor}  
\begin{lstlisting}
    predicate immediateSuccessor { 
      doc /* SFS:  true when this nearly meets o */
      in o : Occurrence;
      nearlyMeets(this,o)
    }
\end{lstlisting}

\pn Use predicate \texttt{timeEnclosedOccurrence} instead of feature \texttt{timeEnclosedOccurrences}.

\ecaption{timeEnclosedOccurrence}  
\begin{lstlisting}
    predicate timeEnclosedOccurrence {
      doc /* SFS:  true when o during self */
      in o : Occurrence;
      during(o,self)
    }
\end{lstlisting}

\pn Use predicate \texttt{timeCoincidentOccurrence} instead of feature \texttt{timeCoincidentOccurrences}.

\ecaption{timeCoincidentOccurrence}  
\begin{lstlisting}
		predicate timeCoincidentOccurrence {
			doc /* SFS:  true when o coincident self */
			in o : Occurrence;
			coincident(o,self)
		}
\end{lstlisting}

\pn Use predicate \texttt{spaceEnclosedOccurrence} instead of feature \texttt{spaceEnclosedOccurrences}.

\ecaption{spaceEnclosedOccurrence}  
\begin{lstlisting}
		predicate spaceEnclosedOccurrence {
			doc /* SFS:  true when self contains o */
			in o : Occurrence;
			regionContainment(location(self),location(o))
		}
\end{lstlisting}

\pn Use predicate \texttt{spaceTimeEnclosedOccurrence} instead of feature \texttt{spaceTimeEnclosedOccurrences}.

\ecaption{spaceTimeEnclosedOccurrence}  
\begin{lstlisting}
		predicate spaceTimeEnclosedOccurrence {
			doc /* SFS:  replaces feature spaceEnclosedOccurrences; true when self contains o */
			in o : Occurrence;
			spaceEnclosedOccurrence(o) and timeEnclosedOccurrence(o)
		}
\end{lstlisting}

\pn Use predicate \texttt{spaceTimeEnclosedPoint} instead of feature \texttt{spaceTimeEnclosedPoints}.

\ecaption{spaceTimeEnclosedPoint}  
\begin{lstlisting}
		predicate spaceTimeEnclosedPoint {
			doc /* SFS: an occurrence with dimension 0 is different than Regions::Point */
			in o : Occurrence;
			spaceTimeEnclosedOccurrence(o) and (o.innerSpaceDimension == 0)
		}
\end{lstlisting}

\pn Predicate \texttt{spaceCoincidentOccurrence} is used by \texttt{spaceTimeCoincidentOccurrence}.

\ecaption{spaceCoincidentOccurrence}  
\begin{lstlisting}
		predicate spaceCoincidentOccurrence {
			doc /* SFS:  true when self and o are space coincident */
			in o : Occurrence;
			regionContainment(location(o),location(self)) and 
			  regionContainment(location(self),location(o))
		}
\end{lstlisting}

\pn Use predicate \texttt{spaceTimeCoincidentOccurrence} instead of feature \texttt{spaceTimeCoincidentOccur\-rences}.

\ecaption{spaceTimeCoincidentOccurrence}  
\begin{lstlisting}
		predicate spaceTimeCoincidentOccurrence {
			doc /* SFS:  true when self and o are space-time coincident */
			in o : Occurrence;
			spaceCoincidentOccurrence(o) and timeCoincidentOccurrence(o)
		}
\end{lstlisting}

\pn Use predicate \texttt{outsideOfOccurrence} instead of feature \texttt{outsideOfOccurrences}.

\ecaption{outsideOfOccurrence}  
\begin{lstlisting}
	  predicate outsideOfOccurrence {
			doc /* SFS:  true when region of self does not overlap region of o */
			in o : Occurrence;
			not regionOverlap(location(self),location(o))
		}
\end{lstlisting}

\pn Use predicate \texttt{justOutsideOfOccurrence} instead of feature \texttt{justOutsideOfOccurrences}.

\ecaption{justOutsideOfOccurrence}  
\begin{lstlisting}
		predicate justOutsideOfOccurrence {
			doc /* SFS:  true when self is just outside of o */
			in o : Occurrence;
			filmConnected(location(self),location(o))
		}
\end{lstlisting}

\pn Use predicate \texttt{matesWithOccurrence} instead of feature \texttt{matesWithOccurrences}.

\ecaption{matesWithOccurrence}  
\begin{lstlisting}
		predicate matesWithOccurrence {
			doc /* SFS:  true when self is just outside of o and o is contained in that */
			in o : Occurrence;
			feature matingSpaceToo: Occurrence[1] subsets self.that;
			justOutsideOfOccurrence(o) and 
			  regionContainment(location(matingSpaceToo),location(o))
		}	
\end{lstlisting}

%\pn Use predicate \texttt{} instead of feature \texttt{}.
%
%\ecaption{}  
%\begin{lstlisting}
%\end{lstlisting}

\subsection{Composite}\label{subsec:composite}

\pn Values of a \texttt{composite} feature, on each instance of the feature's domain, cannot exist after the
featuring instance ceases to exist.  As such, \texttt{composite} features nicely express containment.
For \texttt{suboccurrences}, Domain expresses both mereological and temporal containment.

\ecaption{suboccurrences}  
\begin{lstlisting}
		composite feature suboccurrences: Occurrence[0..*] subsets occurrences {
			doc /* Composite suboccurrences of this Occurrence. */

      language "Domain" /* forall s in suboccurrences are during(life(s),life(self))
			  * and PartOf(s,self) */
        			 
			feature redefines localClock default (that as Occurrence).localClock { }
			 
			feature redefines incomingTransferSort default (that as Occurrence).incomingTransferSort;
		}
\end{lstlisting}

\subsection{Portion}\label{subsec:portion}

\pn Values of a \texttt{portion} feature cannot exist without the whole, because they are the ``same
thing" as the whole.  Although seemingly innocuous, \texttt{portion} features force inheritance upon their values
to be the ``same thing".  

\pn Features \texttt{startShot}, \texttt{endShot}, and \texttt{middleTimeSlice} are appropriate portions; many other features are not.
Space slices and space shots are portions, when sub-components (as suboccurrences) are certainly not the  ``same
thing".

\pn Features \texttt{portions} and \texttt{portionOf} have been removed, as well as inappropriate uses of \texttt{portion}. 

\subsubsection{No Temporal Parts}\label{subsec:ntp}

\pn The metaphysical choice of endurance instead of perdurance (\ref{sec:pervend}) precludes temporal parts.
M0 time is expressed as ``worlds" of \logicid~Logic instead.  Therefore \texttt{timeSlices} and \texttt{snapShots} have
been removed. 

\subsubsection{Space Slices}\label{subsec:spaceslices}

Space slices and shots are no longer portions.

\comment{Conrad Bock has recognized that attempting to combine 3D, 2D, 1D, and 0D (points) has serious problems,
and is writing a paper suggesting resolutions.  The Region library may be helpful.}
%\pn Some features of \texttt{Occurrence} represent model-level (M1) timing relationships.  Such features will be given temporal semantics using \texttt{@Assert}.
%Other features of \texttt{Occurrence} represent execution-level (M0) timing relationships which will be replaced by predicates.
%
%\pn Feature \texttt{sameLifeOccurrences} uses the \texttt{life} function defined in \ref{eq:df-bl.lifetime} to assert that the lifetimes are equal.
%\texttt{sameLifeOccurrences}  are occurrences which are required to be coincident by design (M1)--not those which happen to be coincident (M0).
%\begin{lstlisting}
%feature sameLifeOccurrences: Occurrence[1..*] 
%    subsets things {
%  doc /* sameLifeOccurrences are coincident */
%  @Assert{f="forall o in sameLifeOccurrences are coincident(life(self),life(o)) and"+
%        "o.openLeft=openLeft and o.openRight=openRight"; }
%}
%\end{lstlisting}
%
%\begin{restatable}{definition}{sameLifeOccurrences}~\blTitle{df-sameLifeOccurrences}~ Define \texttt{sameLifeOccurrences}.  
%\begin{align}
%\forall o \in \texttt{sameLifeOccurrences}~\vert~\texttt{coincident}(\texttt{life}(\texttt{self}),\texttt{life}(o))~\wedge \notag \\
%~~\texttt{o.openLeft}=\texttt{self.openLeft}~\wedge~\texttt{o.openRight}=\texttt{self.openRight} 
%\label{eq:df-sameLifeOccurrences}\tag{df-sameLifeOccurrences}
%\end{align}
%\end{restatable}
%
%By definition of \texttt{composite} in \ref{par:cnp}, all \texttt{suboccurrences} have lifetimes during the lifetime of \texttt{self}.
%No annotation is necessary for this temporal inclusion.
%\begin{lstlisting}
%composite feature suboccurrences: Occurrence[0..*] 
%  subsets occurrences {
%	doc /* Composite suboccurrences of this Occurrence. */	 
%  @Assert{f="forall s in suboccurrences are during(life(s),life(self))";}
%  feature redefines localClock default (that as Occurrence).localClock {  }		 
%}
%\end{lstlisting}
%
%\begin{restatable}{definition}{suboccurrences}~\blTitle{df-suboccurrences}~ Define \texttt{suboccurrences}.  
%\begin{align}
%\forall s \in \texttt{suboccurrences}~\vert~\texttt{during}(\texttt{life}(s),\texttt{life}(\texttt{self}))\label{eq:df-suboccurrences}\tag{df-suboccurrences}
%\end{align}
%\end{restatable}
%
%\pn \texttt{predecessors} seems like another M0 object. 
%\begin{lstlisting}
%feature predecessors: Occurrence[0..*]  {
%  @Assert{f="forall p in predecessors are precedes(life(p),life(self)":}
%}
%\end{lstlisting}
%
%\pn Instead, define predecessor\index{predecessor} as a function, invoked as needed:
%\begin{lstlisting}
%predicate predecessor  :> BooleanEvaluation {
%  doc /* SFS:  replaces feature predecessor; true when o precedes self */
%  in o : Occurrence;
%  Time::precedes(o,self)
%}
%\end{lstlisting}
%\begin{restatable}{definition}{predecessor}~\blTitle{df-predecessor}~ Define \texttt{predecessor}.  
%\begin{align}
%\texttt{predecessor}(o)~\equiv  \texttt{precedes}(\texttt{life}(o),\texttt{life}(\texttt{self}))~\label{eq:df-predecessor}\tag{df-predecessor}
%\end{align}
%\end{restatable}
%
%
%\pn \texttt{successors} seems like another M0 object. 
%
%\begin{lstlisting}
%feature successors: Occurrence[0..*] {
%  @Assert{f="forall s in successors are precedes(life(self),life(s)":}
%}
%\end{lstlisting}
%
%\pn Instead, define \texttt{successor}\index{successor} as a predicate, invoked as needed:
%\begin{lstlisting}
%predicate successor  :> BooleanEvaluation {
%  doc /* SFS:  replaces feature successors; true when self precedes o */
%  in o : Occurrence;
%  Time::precedes(self,o)
%}
%\end{lstlisting}
%\begin{restatable}{definition}{successor}~\blTitle{df-successor}~ Define \texttt{successor}.  
%\begin{align}
%\texttt{successor}(o)~\equiv~ \texttt{precedes}(\texttt{life}(\texttt{self}),\texttt{life}(o))~\}\label{eq:df-successor}\tag{df-successor}
%\end{align}
%\end{restatable}
%
%\pn An immediate predecessor ceases to exist an infinitesimal time, \texttt{delta} before this occurrence comes into existence.
%No other occurrence may occur between them; it nearly meets.
%\begin{lstlisting}
%feature immediatePredecessors: Occurrence[0..*] {
%	doc /* Occurrences that end just before this occurrence starts, with no
%	* possibility of other occurrences happening in the time between them.
%	*/
%	@Assert{f="forall i in immediatePredecessors are nearlyMeets(i,self)";}
%}
%\end{lstlisting}
%
%\begin{restatable}{definition}{impred}~\blTitle{df-impred}~ Define \texttt{immediatePredecessors}.  
%\begin{align}
%\forall i \in \texttt{immediatePredecessors}~\vert~\texttt{nearlyMeets}(\texttt{life}(i),\texttt{life}(\texttt{self}))\label{eq:df-impred}\tag{df-impred}
%\end{align}
%\end{restatable}
%
%
%\pn Immediate successors are the converse of immediate predecessors.  The immediate successors of self is not the inverse of the immediate predecessors of self.
%\begin{lstlisting}
%feature immediateSuccessors: Occurrence[0..*] {
%	doc
%	/* Occurrences that start just after this occurrence ends, with no
%	* possibility of other occurrences happening in the time between them.
%	*/
%	@Assert{f="forall i in immediateSuccessors are nearlyMeets(self,i)";}
%}\end{lstlisting}
%
%\begin{restatable}{definition}{imsucc}~\blTitle{df-imsucc}~ Define \texttt{immediateSuccessors}.  
%\begin{align}
%\forall i \in \texttt{immediateSuccessors}~\vert~\texttt{nearlyMeets}(\texttt{life}(\texttt{self}),\texttt{life}(i))\label{eq:df-imsucc}\tag{df-imsucc}
%\end{align}
%\end{restatable}
%
%\pn Time-enclosed occurrences are those which just happen during self (M0), unlike suboccurrences which must happen during self (M1).  %The other features seem unnecessary, and possibly incorrect.
%\begin{lstlisting}
%feature timeEnclosedOccurrences: Occurrence[1..*] subsets occurrences {
%  doc
%	/*
%	 * Occurrences that start no earlier than and end no later than
%	 * this occurrence, including at least this occurrence.
%	 */
%  @Assert{f="forall teo in timeEnclosedOccurrences are during(teo,self)";}
%}
%\end{lstlisting}
%\pn Instead, define \texttt{timeEnclosedOccurrence} as a predicate, invoked as needed:
%\begin{lstlisting}
%predicate timeEnclosedOccurrence  :> BooleanEvaluation {
%  doc /* SFS:  replaces feature timeEnclosedOccurrences; true when o during self */
%  in o : Occurrence;
%  Time::during(o,self)
%}
%\end{lstlisting}
%
%By declaring the feature as \bemph{portion}, portions is already constrained to happen during self.
%\begin{lstlisting}
%portion feature all portions: Occurrence[1..*] subsets spaceTimeEnclosedOccurrences {
%  @Assert{f="forall p in portions are during(p,self)";} }
%\end{lstlisting}

%\comment{I have a deep philosophical disagreement with defining time instants (snapshots) as interval with zero duration.
%\texttt{snapshots} are defined as occurrences where \texttt{startShot[1]=endShot[1]}.  But then \texttt{startShot} subsets \texttt{snapshots}.
%So startShot.startShot = startShot.endShot, and startShot.startShot.startShor = startShot.endShot.startShort, and so on.  }

%\comment{Instead define type \texttt{Time} typed by \texttt{Real} in \texttt{ScalarValues}.}

%\texttt{startShot} is the birth of the classifier.
%\begin{lstlisting}
%portion feature startShot: Instant[1] {
%  language "assertion" /* timeof(startShot) = birth(self) */
%\end{lstlisting}
%
%\texttt{endShot} is the death of the classifier.
%\begin{lstlisting}
%portion feature endShot: Instant[1] {
%  language "assertion" /* timeof(endShot) = death(self) */
%\end{lstlisting}
%
%\texttt{middleTimeSlice} is the open interval between \texttt{startShot} and \texttt{endShot}--if  they're different.
%However, \texttt{Occurrence} is defined as a closed interval.  Here, \texttt{delta} is useful. 
%\begin{lstlisting}
%portion feature middleTimeSlice: Occurrence[0..1] {
%  language "assertion" 
%  /* middleTimeSlice.startShot = birth(self)+delta and 
%  * middleTimeSlice.endShot+delta = death(self) 
%  */
%\end{lstlisting}

\subsection{Associations}\label{subsec:associations}

\pn The \texttt{Occurrences} library defines many associations, most of which include \texttt{all} indicating sufficiency.
Does the \texttt{all} mean that every possible association of that type exists?  This seems wrong.

\pn \texttt{SelfSameLifeLink} is convoluted because it tries to treat DataValues and Occurrences in a uniform way. 
\texttt{Links::SelfLink} should suffice.

\pn \texttt{Within}, \texttt{WithinBoth}, \texttt{PortionOf}, \texttt{TimeSliceOf}, \texttt{SnapshotOf}, and  \texttt{Without} are unnecessary.

\pn \texttt{HappensDuring}, \texttt{HappensWhile}, \texttt{HappensBefore}, and \texttt{HappensJustBefore} lose their  \texttt{all} and gain temporal semantics.

\ecaption{HappensDuring}  
\begin{lstlisting}
	assoc HappensDuring specializes HappensLink {
		doc /*
		* HappensDuring asserts that the shorterOccurrence happens during the longerOccurrence. */
		language "Domain" /* during(shorterOccurrence, longerOccurrence) */
		end feature shorterOccurrence redefines sourceOccurrence;  
		end feature longerOccurrence redefines targetOccurrence;
	}
\end{lstlisting}

\ecaption{HappensWhile}  
\begin{lstlisting}
	assoc  HappensWhile specializes HappensLink {
		doc
		/*
		 * HappensWhile asserts that two occurrences happen during each other, that is, they
		 * each start at the same time and end at the same time.
		 */
		language "Domain" /* coincident(thisOccurrence, thatOccurrence) */
		end feature thisOccurrence redefines sourceOccurrence;  
		end feature thatOccurrence redefines targetOccurrence;
	}
\end{lstlisting}

\ecaption{HappensBefore}  
\begin{lstlisting}
	assoc HappensBefore specializes HappensLink 
	  {
		language "Domain" /* precedes(earlierOccurrence,laterOccurrence) */
		end feature earlierOccurrence: Occurrence redefines sourceOccurrence; 
		end feature laterOccurrence: Occurrence redefines targetOccurrence; 
	}
\end{lstlisting}

\ecaption{HappensJustBefore}  
\begin{lstlisting}
	assoc HappensJustBefore specializes HappensBefore {
		language "Domain" /* nearlyMeets(earlierOccurrence,laterOccurrence) */
		end feature redefines earlierOccurrence: Occurrence
		  crosses laterOccurrence.immediatePredecessors;
		end feature redefines laterOccurrence: Occurrence
		  crosses earlierOccurrence.immediateSuccessors;
	}
\end{lstlisting}

\pn \texttt{InsideOf}, \texttt{OutsideOf}, \texttt{JustOutsideOf}, \texttt{MatesWith}, and \texttt{InnerSpaceOf} 
lose the \texttt{all} and gain regional semantics.

\ecaption{InsideOf}  
\begin{lstlisting}
	assoc InsideOf specializes SpaceLink {
		language "Domain" /* regionContainment(location(smallerSpace), location(largerSpace)) */
		end feature smallerSpace: Occurrence redefines sourceOccurrence;  
		end feature largerSpace: Occurrence redefines targetOccurrence;
	}
\end{lstlisting}

\ecaption{OutsideOf}  
\begin{lstlisting}
	assoc OutsideOf specializes SpaceLink 
	  {
		language "Domain" 
		  /* not regionOverlap(location(separateSpaceToo), location(separateSpace)) */
		end feature separateSpaceToo: Occurrence redefines sourceOccurrence; 
		end feature separateSpace: Occurrence redefines targetOccurrence; 
	}
\end{lstlisting}

\ecaption{JustOutsideOf}  
\begin{lstlisting}
	assoc JustOutsideOf specializes OutsideOf {
		doc
		/*
		 * JustOutsideOf is an OutsideOf asserting that two occurrences have some space slices 
		 * with no space between them.
		 */
		language "Domain" /* filmConnected(separateSpaceToo, separateSpace) */
		end feature redefines separateSpaceToo: Occurrence;
		end feature redefines separateSpace: Occurrence;
	}
\end{lstlisting}

\comment{not sure how MatesWith differs from JustOutsideOf}
\ecaption{MatesWith}  
\begin{lstlisting}
	assoc MatesWith specializes JustOutsideOf {
		doc
		/*
		 * MatesWith is an OutsideOf asserting that two occurrences have no space between them.
		 */
		end feature matingSpaceToo: Occurrence redefines separateSpaceToo;
		end feature matingSpace: Occurrence redefines separateSpace;
	}
\end{lstlisting}

\ecaption{InnerSpaceOf}  
\begin{lstlisting}
	assoc InnerSpaceOf specializes OutsideOf {
		doc
		/*
		 * InnerSpaceOf is an OutsideOf asserting that the space surrounded by an inner space boundary
		 * of one occurrence (outer space) is completely occupied by another occurrence (inner space).
		 */
		language "Domain" /* regionFilm(location(innerSpace)) = regionSurface(location(outerSpace))*/
		end feature outerSpace: Occurrence redefines separateSpaceToo;
		end feature innerSpace: Occurrence redefines separateSpace; 
	}
\end{lstlisting}



